%!TEX TS-program = xelatex
\documentclass[12pt, a4paper, oneside]{article}

% пакеты для математики
\usepackage{amsmath,amsfonts,amssymb,amsthm,mathtools}  
\mathtoolsset{showonlyrefs=true}  % Показывать номера только у тех формул, на которые есть \eqref{} в тексте.

\usepackage[british,russian]{babel} % выбор языка для документа
\usepackage[utf8]{inputenc}          % utf8 кодировка

% Основные шрифты 
\usepackage{fontspec}         
\setmainfont{Linux Libertine O}  % задаёт основной шрифт документа

% Математические шрифты 
\usepackage{unicode-math}     
\setmathfont[math-style=upright]{Neo Euler} 

%%%%%%%%%% Работа с картинками и таблицами %%%%%%%%%%
\usepackage{graphicx} % Для вставки рисунков                
\usepackage{graphics}
\graphicspath{{images/}{pictures/}}   % папки с картинками

\usepackage[figurename=Картинка]{caption}

\usepackage{wrapfig}    % обтекание рисунков и таблиц текстом

\usepackage{booktabs}   % таблицы как в годных книгах
\usepackage{tabularx}   % новые типы колонок
\usepackage{tabulary}   % и ещё новые типы колонок
\usepackage{float}      % возможность позиционировать объекты в нужном месте
\renewcommand{\arraystretch}{1.2}  % больше расстояние между строками


%%%%%%%%%% Графики и рисование %%%%%%%%%%
\usepackage{tikz, pgfplots}  % языки для графики
%\pgfplotsset{compat=1.16}

\usepackage{todonotes} % для вставки в документ заметок о том, что осталось сделать
% \todo{Здесь надо коэффициенты исправить}
% \missingfigure{Здесь будет Последний день Помпеи}
% \listoftodos --- печатает все поставленные \todo'шки


%%%%%%%%%% Внешний вид страницы %%%%%%%%%%

\usepackage[paper=a4paper, top=20mm, bottom=15mm,left=20mm,right=15mm]{geometry}
\usepackage{indentfirst}    % установка отступа в первом абзаце главы

\usepackage{setspace}
\setstretch{1.15}  % межстрочный интервал
\setlength{\parskip}{4mm}   % Расстояние между абзацами
% Разные длины в LaTeX: https://en.wikibooks.org/wiki/LaTeX/Lengths

% свешиваем пунктуацию
% теперь знаки пунктуации могут вылезать за правую границу текста, при этом текст выглядит ровнее
\usepackage{microtype}

% \flushbottom                            % Эта команда заставляет LaTeX чуть растягивать строки, чтобы получить идеально прямоугольную страницу
\righthyphenmin=2                       % Разрешение переноса двух и более символов
\widowpenalty=300                     % Небольшое наказание за вдовствующую строку (одна строка абзаца на этой странице, остальное --- на следующей)
\clubpenalty=3000                     % Приличное наказание за сиротствующую строку (омерзительно висящая одинокая строка в начале страницы)
\tolerance=10000     % Ещё какое-то наказание.

% мои цвета https://www.artlebedev.ru/colors/
\definecolor{titleblue}{rgb}{0.2,0.4,0.6} 
\definecolor{blue}{rgb}{0.2,0.4,0.6} 
\definecolor{red}{rgb}{1,0,0.2} 
\definecolor{green}{rgb}{0,0.6,0} 
\definecolor{purp}{rgb}{0.4,0,0.8} 

% цвета из geogebra 
\definecolor{litebrown}{rgb}{0.6,0.2,0}
\definecolor{darkbrown}{rgb}{0.75,0.75,0.75}

% Гиперссылки
\usepackage{xcolor}   % разные цвета

\usepackage{hyperref}
\hypersetup{
	unicode=true,           % позволяет использовать юникодные символы
	colorlinks=true,       	% true - цветные ссылки
	urlcolor=blue,          % цвет ссылки на url
	linkcolor=black,          % внутренние ссылки
	citecolor=green,        % на библиографию
	breaklinks              % если ссылка не умещается в одну строку, разбивать её на две части?
}

% меняю оформление секций 
\usepackage{titlesec}
\usepackage{sectsty}

% меняю цвет на синий
\sectionfont{\color{titleblue}}
\subsectionfont{\color{titleblue}}

% выбрасываю нумерацию страниц и колонтитулы 
\pagestyle{empty}

% синие круглые бульпоинты в списках itemize 
\usepackage{enumitem}

\definecolor{itemizeblue}{rgb}{0, 0.45, 0.70}

\newcommand*{\MyPoint}{\tikz \draw [baseline, fill=itemizeblue, draw=blue] circle (2.5pt);}
\renewcommand{\labelitemi}{\MyPoint}

\AddEnumerateCounter{\asbuk}{\@asbuk}{\cyrm}
\renewcommand{\theenumi}{\asbuk{enumi}}

% расстояние в списках
\setlist[itemize]{parsep=0.4em,itemsep=0em,topsep=0ex}
\setlist[enumerate]{parsep=0.4em,itemsep=0em,topsep=0ex}

% эпиграфы
\usepackage{epigraph}
\setlength\epigraphwidth{.4\textwidth}
\setlength\epigraphrule{0pt}

%%%%%%%%%% Свои команды %%%%%%%%%%

% Математические операторы первой необходимости:
\DeclareMathOperator{\sgn}{sign}
\DeclareMathOperator*{\argmin}{arg\,min}
\DeclareMathOperator*{\argmax}{arg\,max}
\DeclareMathOperator{\Cov}{Cov}
\DeclareMathOperator{\Var}{Var}
\DeclareMathOperator{\Corr}{Corr}
\DeclareMathOperator{\E}{\mathop{E}}
\DeclareMathOperator{\Med}{Med}
\DeclareMathOperator{\Mod}{Mod}
\DeclareMathOperator*{\plim}{plim}

% команды пореже
\newcommand{\const}{\mathrm{const}}  % const прямым начертанием
\newcommand{\iid}{\sim i.\,i.\,d.}  % ну вы поняли...
\newcommand{\fr}[2]{\ensuremath{^{#1}/_{#2}}}   % особая дробь
\newcommand{\ind}[1]{\mathbbm{1}_{\{#1\}}} % Индикатор события
\newcommand{\dx}[1]{\,\mathrm{d}#1} % для интеграла: маленький отступ и прямая d

% одеваем шапки на частые штуки
\def \hb{\hat{\beta}}
\def \hs{\hat{s}}
\def \hy{\hat{y}}
\def \hY{\hat{Y}}
\def \he{\hat{\varepsilon}}
\def \hVar{\widehat{\Var}}
\def \hCorr{\widehat{\Corr}}
\def \hCov{\widehat{\Cov}}

% Греческие буквы
\def \a{\alpha}
\def \b{\beta}
\def \t{\tau}
\def \dt{\delta}
\def \e{\varepsilon}
\def \ga{\gamma}
\def \kp{\varkappa}
\def \la{\lambda}
\def \sg{\sigma}
\def \tt{\theta}
\def \Dt{\Delta}
\def \La{\Lambda}
\def \Sg{\Sigma}
\def \Tt{\Theta}
\def \Om{\Omega}
\def \om{\omega}

% Готика
\def \mA{\mathcal{A}}
\def \mB{\mathcal{B}}
\def \mC{\mathcal{C}}
\def \mE{\mathcal{E}}
\def \mF{\mathcal{F}}
\def \mH{\mathcal{H}}
\def \mL{\mathcal{L}}
\def \mN{\mathcal{N}}
\def \mU{\mathcal{U}}
\def \mV{\mathcal{V}}
\def \mW{\mathcal{W}}

% Жирные буквы
\def \mbb{\mathbb}
\def \RR{\mbb R}
\def \NN{\mbb N}
\def \ZZ{\mbb Z}
\def \PP{\mbb{P}}
\def \QQ{\mbb Q}

%%%%%%%%%% Теоремы %%%%%%%%%%
\theoremstyle{plain} % Это стиль по умолчанию.  Есть другие стили.
\newtheorem{theorem}{Теорема}[section]
\newtheorem{result}{Следствие}[theorem]
% счётчик подчиняется теоремному, нумерация идёт по главам согласованно между собой

% убирает курсив и что-то еще наверное делает ;)
\theoremstyle{definition}         
\newtheorem*{definition}{Определение}  % нумерация не идёт вообще


%%%%%%%%%% Задачки и решения %%%%%%%%%%
\usepackage{etoolbox}    % логические операторы для своих макросов
\usepackage{environ}
\newtoggle{lecture}

\newcounter{probNum}[section]  % счётчик для упражнений 
\NewEnviron{problem}[1]{%
	\refstepcounter{probNum}% увеличели номер на 1 
	{\noindent \textbf{\large \color{titleblue} Упражнение~\theprobNum~#1}  \\ \\ \BODY}
	{}%
}

% Окружение, чтобы можно было убирать решения из pdf
\NewEnviron{sol}{%
	\iftoggle{lecture}
	{\noindent \textbf{\large Решение:} \\ \\ \BODY}
	{}%
}

% выделение по тексту важных вещей
\newcommand{\indef}[1]{\textbf{ \color{green} #1}} 

% разные дополнения для картинок
\usetikzlibrary{arrows.meta}
\usepackage{varwidth}

% рисование крестов
% https://tex.stackexchange.com/questions/123760/draw-crosses-in-tikz
\tikzset{
	cross/.pic = {
		\draw[line width=2.pt, rotate = 45] (-#1,0) -- (#1,0);
		\draw[line width=2.pt, rotate = 45] (0,-#1) -- (0, #1);
	}
}

\usepackage[normalem]{ulem}  % для зачекивания текста
\usepackage{tikz, pgfplots}  % язык для рисования графики из latex'a

\usepackage{todonotes} % для вставки в документ заметок о том, что осталось сделать
% \todo{Здесь надо коэффициенты исправить}
% \missingfigure{Здесь будет Последний день Помпеи}
% \listoftodos --- печатает все поставленные \todo'шки

\DeclareMathOperator{\tr}{tr}

\title{%
	Тятя! Тятя! Наши сети притащили мертвеца! \\
	\large Листочек с задачками №3: Матричное дифФфФфФфириенцирование\footnote{Часть задач взята из \href{https://github.com/bdemeshev/mlearn_pro/blob/master/mlearn_pro.pdf}{прототипа задачника по ML  Бориса Демешева}, часть из \href{https://github.com/esokolov/ml-course-msu}{конспектов по ML Жени Соколова}}
}
\date{РАНХ \\ осень 2020}
\author{\url{https://github.com/FUlyankin/neural_nets_prob} }



\begin{document}
	
	\maketitle
	
\epigraph{ $\bigg(\text{ \includegraphics[scale=0.08]{tree1}} \bigg)^T = $  \includegraphics[scale=0.08]{tree2}}{<<Джек и бобовый стебель>> (1890)}


%%%-------------------------------------------
\begin{problem}{}	
	В этой задачке нужно  просто найти немного производных: 
	\begin{enumerate}	
		
		\item $f(x) = a^T x$, где  $a$  и $x$ векторы размера $1 \times n$ 
		
		\item $f(x) = x^T A x$, где $x$ вектор размера $1 \times n$, $A$ матрица размера $n \times n$
		
		\item $f(x) = \ln(x^T A x)$, где $x$ вектор размера $1 \times n$, $A$ матрица размера $n \times n$
		
		\item $f(x) = a^TXAXa$, где $x$ вектор размера $1 \times n$, $A$ матрица размера $n \times n$
		
		\item $f(x) = x x^T x$,  где $x$ вектор размера $1 \times n$
		
		\item $f(X) = X^{-1}$, где матрица $X$ размера $n \times n$
		
		\item  $f(X) = \det X$, где матрица $X$ размера $n \times n$
	\end{enumerate}
\end{problem} 


%%%-------------------------------------------
\begin{problem}{}
	В этой задачке нужно  просто найти много разных производных: 
	\begin{enumerate}
		
		
		\item $f(X) = \tr(AXB)$, где матрица $A$ размера $p \times m$, матрица $B$ размера $n \times p$, матрица $X$ размера $m \times n$. 
		
		\item $f(X) = \tr(AX^TX)$, где матрица $A$ размера $n \times n$, матрица $X$ размера $m \times n$. 
		
		\item $f(X) = \ln \det X$
		
		\item $f(X) = \ln AX^{-1}B$  % убрать эту задачу 
		
		\item $f(X) = \tr(AX^TXBX^{-T})$
		
		\item $f(X) = \ln \det(X^TAX)$
		
		\item $f(x) = x^TAb$, где матрица $A$ размера $n \times n$, вектора $x$ и $b$ размера $n \times 1$. 
		
		\item $f(A) = x^TAb$. 
	\end{enumerate}
\end{problem}

\begin{problem}{}
	Рассмотрим задачу линейной регресии
	\[
	Q(w) = (y - Xw)^T(y - Xw) \to \min_{w}.
	\]
	
	\begin{enumerate}
		\item Найдите $dQ(w)$, выведите формулу для оптимального $w$.
		\item Как выглядит шаг градиентного спуска в матричном виде?
		\item Найдите $d^2Q(w)$. Убедитесь, что мы действительно в точке минимума. 
	\end{enumerate}
\end{problem}

\begin{problem}{}
	В случае Ridge-регрессии минимизируется функция
	\[
	Q(w) = (y - Xw)^T(y - Xw) + \lambda w^T w,
	\]
	где $\lambda$ — положительный параметр, штрафующий функцию за слишком большие значения $w$.
	
	\begin{enumerate}
		\item Найдите $dQ(w)$, выведите формулу для оптимального $w$.
		\item Как выглядит шаг градиентного спуска в матричном виде?
		\item Найдите $d^2Q(w)$. Убедитесь, что мы действительно в точке минимума. 
	\end{enumerate}
	
	В случае Lasso-регрессии мы имеем дело с функцией
	\[
	Q(w) = (y - Xw)^T(y - Xw) + \lambda |w|,
	\]
	
	\begin{enumerate}
		\item Найдите $dQ(w)$, выведите формулу для оптимального $w$.
		\item Как выглядит шаг градиентного спуска в матричном виде?
	\end{enumerate}
\end{problem}

\begin{problem}{}
	
	Пусть $x_i$ — вектор-столбец $k\times 1$, $y_i$ — скаляр, равный $+1$ или $-1$, $w$ — вектор-столбец размера $k\times 1$. Рассмотрим функцию
	\[
	Q(w) = \sum_{i=1}^n \ln (1 + \exp(-y_ix_i^Tw)) + \lambda w^T w
	\]
	
	\begin{enumerate}
		\item Найдите $dQ$;
		\item Найдите вектор-столбец $\nabla Q$.
	\end{enumerate}
\end{problem}


\begin{problem}{}
	Упражняемся в матричном методе максимального правдоподобия.  Допустим, что векторы $X_1, \ldots, X_m$ выбраны из многомерного нормального распределения с неизвестными вектором средних $\mu$ и ковариационной матрицей $\Sigma$. В этом задании нужно найти оценки максимального правдоподобия для $\hat \mu$ и $\hat \Sigma$.  Обратите внимание, что выборкой здесь будет не $x_1, \ldots, x_m$, а 
	\begin{equation*}
	\begin{pmatrix}
	x_{11}, \ldots, x_{m1} \\
	\ldots  \\ 
	x_{1n}, \ldots, x_{mn}
	\end{pmatrix}
	\end{equation*}
\end{problem}


\begin{problem}{}
	Найдите симметричную матрицу $X$ наиболее близкую к матрице $A$ по норме Фробениуса, $\sum_{i,j} (x_{ij} - a_{ij})^2$. Тут мы просто из каждого элемента вычитаем каждый и смотрим на сумму квадратов таких разностей. 
	
	То есть решите задачку условной матричной минимизации 
	
	\begin{equation*}
	\begin{cases}
	& ||X - A||^2 \to \min_{A}  \\
	& X^T = X
	\end{cases}
	\end{equation*}
	
	\textbf{Hint:} Надо будет выписать Лагранджиан.  А ещё пригодится тот факт, что $\sum_{i,j} (x_{ij} - a_{ij})^2 = ||X-A||^2 =  \tr((X-A)^T (X-A))$. То, что это так мы доказали на семинаре :) Вспоминайте! 
\end{problem}



\end{document}